The Kadomtsev--Petviashvili (KP) hierarchy and
the Toda lattice hierarchy are two important integrable hierarchies. In this
paper, we will show that the tau-function of an arbitrary solution to the Toda lattice hierarchy
gives an infinite family of tau-functions of the KP hierarchy.
This is achieved by combining two different results:
a formula~\cite{Zhou-Emergent} of the second-named author on the $n$-point function for a KP tau-function in the big cell,
and a work~\cite{Y} of the first-named author on the $n$-point function for a Toda tau-function.

The KP hierarchy is an infinite family of equations with infinitely many unknown functions,
which can be written using the Lax pair formalism as follows:
\beq\label{KPhierarchy30}
\frac{\p L_{\rm KP}}{\p T_k} = \bigl[\bigl(L_{\rm KP}^k\bigr)_+, L_{\rm KP}\bigr], \qquad k\ge1,
\eeq
where
\[
L_{\rm KP} = \p + \sum_{j\ge1} u_j \p^{-j}, \qquad \p=\p/\p X,
\]
is a pseudodifferential operator, called the {\it Lax operator}.
For details about pseudodifferential operators and their operations see for example~\cite{Dickey}.
The independent variables $T_1,T_2,T_3,\dots$ are called {\it times}.
Since $\p L_{\rm KP}/\p T_1 = \p L_{\rm KP}/\p X$, we identify $T_1$ with $X$.

Let ${\bf T}=(T_1,T_2,\dots)$ denote the infinite vector of KP times.
It is well known (see, e.g.,~\cite{Dickey}) that an arbitrary solution $(u_1({\bf T}), u_2({\bf T}),\dots)$
to the KP hierarchy~\eqref{KPhierarchy30} can be compactly represented by a single function $\tau=\tau({\bf T})$
called the {\it tau-function} as
\[
u_1({\bf T}) = \frac{\p^2 \log \tau}{\p T_1^2}, \qquad u_2({\bf T}) = \frac12 \frac{\p^2 \log \tau}{\p T_1 \p T_2} - \frac12 \frac{\p^3 \log \tau}{\p T_1^3}, \qquad \dots,
\]
and the tau-function~$\tau$ satisfies the {\it Hirota bilinear identities} given by
\beq\label{bilinearidH}
{\rm res}_{\lambda=\infty} \tau\bigl({\bf T} - \bigl[\lambda^{-1}\bigr]\bigr) \tau\bigl({\bf T'} + \bigl[\lambda^{-1}\bigr]\bigr) {\rm e}^{\sum_{k\ge1} (T_k-T_k') \lambda^k} {\rm d}\lambda=0.
\eeq
Here \smash{$\bigl[\lambda^{-1}\bigr]:=\bigl(\frac1\lambda, \frac1{2\lambda^2}, \frac 1{3\lambda^3}, \dots\bigr)$}.
Equivalence between~\eqref{bilinearidH} and~\eqref{KPhierarchy30} is a standard result in the theory of integrable systems; see, for example,~\cite{Dickey}.
Equation~\eqref{bilinearidH} itself gives the defining equations for a KP tau-function (namely, we do not
 need to start with a solution $(u_1({\bf T}), u_2({\bf T}),\dots)$ to the KP hierarchy~\eqref{KPhierarchy30}).

Let $\Lambda\colon f(x)\mapsto f(x+\e)$ be the shift operator and
\begin{gather*}% \label{eqn:Lax-Toda}
L := \Lambda + v(x) + w(x) \Lambda^{-1}
\end{gather*}
the {\it Lax operator}. Here $\e$ is a parameter.
The {\it Toda lattice hierarchy} (also known as the 1d Toda chain or 1-Toda hierarchy)
can be defined using the Lax pair formalism as
\beq\label{TLH}
\e\frac{\p L}{\p m_i} = \frac1{(i+1)!} \bigl[\bigl(L^{i+1}\bigr)_+,L\bigr], \qquad i\geq0.
\eeq
Here, for a difference operator~$P$ written in the form $P=\sum_{k\in \ZZ} P_k \Lambda^k$,
$P_+$ is defined as $\sum_{k\geq0} P_k \Lambda^k$. We also denote ${\bf m}=(m_0,m_1,m_2,\dots)$.

It is known~\cite{UT} that an arbitrary tau-function $\tau(N,{\bf x},{\bf y})$ of the 2-Toda hierarchy~\cite{UT}
is a~KP~tau-function with respect to either ${\bf x}$ or~${\bf y}$ as KP times (fixing the others
as parameters). The~following Theorem~\ref{thm:toda-kp},
which could essentially be deduced from~\cite{UT} (see Remark~\ref{remark12} below), gives a similar statement for the 1-Toda hierarchy. Let $(v(x,{\bf m};\e), w(x,{\bf m};\e))$ be an arbitrary power-series-in-${\bf m}$ solution
to the Toda lattice hierarchy
with coefficients being in a ring $V$ of functions of~$x$ closed under shifting~$x$
by $\pm \e$ (e.g., $V$ could be $\mathbb{C}[x,\e]$),
and $\tau(x,{\bf m};\e)$ the tau-function~\cite{DY} of this solution.
%Our first result of this paper is given by the following
\begin{Theorem} \label{thm:toda-kp}
The tau-function $\tau(x,{\bf m};\e)$ for the Toda lattice hierarchy
is a tau-function of the KP hierarchy for any $x$ and $\e$,
where ${\bf m}$ and the KP times ${\bf T}$ are related by
\be\label{tT30}
m_i = (i+1)! \, \e \, T_{i+1},\qquad i\ge0.
\ee
\end{Theorem}
\begin{Remark}\label{remark12}
Theorem~\ref{thm:toda-kp} should be known to experts, and could be deduced
from a claim~\cite[p.~30]{UT} (cf.\ \cite{KMMM}) that the Lax representation of the
Toda lattice hierarchy can be obtained from
a reduction of the 2-Toda hierarchy with
the reduction condition
\[
\p/\p x_i = \p /\p y_i, \qquad i\ge1,
\]
up to a quadratic function
of ${\bf x}$, ${\bf y}$. Here, ${\bf x}=(x_1,x_2,\dots)$, ${\bf y}=(y_1,y_2,\dots)$.
(Our formula~\eqref{dpsiabtoda} below should shed light on the Toda chain discussions in~\cite{KMMM}.)
We would like to thank A.~Alexandrov for bringing our attention to the reduction arguments of~\cite{UT}.
We also point out the following well-known fact: the first nontrivial Hirota bilinear equation
of the 2-Toda hierarchy under the reduction condition coincides with that of the 1-Toda hierarchy. Thus, an alternative way to
show the equivalence between the 1-Toda hierarchy and the reduction of the 2-Toda hierarchy is
to compare the Hirota bilinear equations of the two hierarchies~\cite{Milanov, UT}, whose detail might deserve
a further study.
In this work we take a different approach. For the KP hierarchy,
a formula for the connected $n$-point functions was derived by the second-named author using the fermionic approach in~\cite{Zhou-Emergent},
whereas for the Toda lattice hierarchy a formula of the same form was derived by the second-named author~\cite{Y}, but using a totally different method:
the matrix-resolvent method. So we present a~different proof combining these two approaches. An advantage of this proof is that it can be generalized
to the extended Toda hierarchy~\cite{CDZ} whose fermionic approach has not yet appeared. See our Theorem~\ref{thm:toda-kp-ext} below. The main content of Section~\ref{sectionTodaext}
is to generalize the matrix-resolvent method to the case of the extended Toda hierarchy.
\end{Remark}
